\documentclass[a4paper]{article}
\usepackage{titlesec}
\usepackage{geometry}
\geometry{hmargin={2cm,2cm}, vmargin={3.5cm,3.5cm}}
\usepackage{graphicx}
\usepackage[table]{xcolor}
\usepackage{hyperref}
\hypersetup{pdfborder = {0 0 0},
colorlinks,
citecolor=black,
filecolor=black,
linkcolor=black,
urlcolor=blue}
\definecolor{logo}{RGB}{20,99,163}
\usepackage[T1]{fontenc}
\usepackage[utf8]{inputenc}
\usepackage[english,italian]{babel}
\usepackage{fancyhdr}
\usepackage{lastpage}
\usepackage{comment}

% creazione comando subsubsubsection
\titleclass{\subsubsubsection}{straight}[\subsection]

\newcounter{subsubsubsection}[subsubsection]
\renewcommand\thesubsubsubsection{\thesubsubsection.\arabic{subsubsubsection}}
\titleformat{\subsubsubsection}
    {\normalfont\normalsize\bfseries}{\thesubsubsubsection}{1em}{}
    \titlespacing*{\subsubsubsection}
    {0pt}{3.25ex plus 1ex minus .2ex}{1.5ex plus .2ex}

\makeatletter
\def\toclevel@subsubsubsection{4}
\def\l@subsubsubsection{\@dottedtocline{4}{7em}{4em}}
\makeatother

\setcounter{secnumdepth}{4}
\setcounter{tocdepth}{4}
%fine creazione comando


%Comandi di impaginazione uguale per tutti i documenti
\pagestyle{fancy}
\lhead{\includegraphics[scale=0.2]{../../../../Images/logo.png}}
%Titolo del documento
\rhead{\large Verbale Esterno del 2021-05-16}
%\rfoot{\thepage}
\cfoot{Pagina \thepage{} di \pageref{LastPage}}
\setlength{\headheight}{45pt}
\renewcommand{\footrulewidth}{0.4pt}
\begin{document}

\begin{titlepage}
    \begin{center}
        \includegraphics{../../../../Images/logo.png}\\
        \vspace{20px}
        \textcolor{logo}{\hrulefill}\\
        \vspace{20px}
        \textbf{\huge\textcolor{logo}{Verbale Esterno del 2021-05-16}}\\
        \vspace{10px}
        \textcolor{logo}{\hrulefill}\\
        \vspace{40px}
        \textbf{\Large Informazioni sul documento}\\
        \vspace{20px}
        \begin{tabular}{p{100px} | p{100px}}
            \textbf{Versione}     & 1.0.0                                                                      \\
            \textbf{Responsabile} & Tito Scutari                                                               \\
            \textbf{Redattore}    & Marco Barbaresco                                                           \\
            \textbf{Verificatore} & Samuele De Simone                                                          \\
            \textbf{Destinatari}  & Prof. Tullio Vardanega \newline Prof. Riccardo Cardin, \newline TechSWEave \\
            \textbf{Stato}        & Approvato                                                                  \\
            \textbf{Uso}          & Esterno                                                                    \\
        \end{tabular}\\
        \vspace{60px}
        \href{mailto:techsweave@gmail.com}{techsweave@gmail.com}\\

    \end{center}
\end{titlepage}

%INDICE
\newpage
\tableofcontents
\newpage

\section{Informazioni riunione}
\begin{itemize}
    \item Data e ora: 2021-05-16 ore 17:00-17:30.
    \item Luogo: Google Meet fornito dai proponenti.
    \item Partecipanti interni: Samuele De Simone, Manuel Varo, Tito Scutari, Nicolò Giaccone, Amedeo Meggiolaro, Marco Barbaresco.
    \item Partecipanti esterni: Milo Ertola.
    \item Segretario: Marco Barbaresco.
\end{itemize}
\section{Argomenti affrontati}
Durante l'incontro sono stati chiesti al proponente chiarimenti e consigli riguardo dettagli progettuali ed implementativi.
\section{Verbale della riunione}
\begin{itemize}
    \item Aggiornamento del proponente sullo stato del progetto;
    \item Richiesti al proponente pareri sul passaggio ad un'architettura a microservizi;
    \item Richiesti al proponente consigli riguardo l'utilizzo di server-side e static-rendering;
    \item Richiesti al proponente consigli sull'utilizzo del framework Bootstrap o altri framework similari utili lato frontend;
    \item Richiesti al proponente consigli su ESLint lato backend;
    \item Richiesti al proponente consigli sulla gestione della sessione dei pagamenti;
    \item Fissata una riunione per il 2021-05-21 in cui saranno presenti entrambi i rappresentanti dell'azienda proponente \textit{Red Babel}.
\end{itemize}
\newpage
\section{Resoconto}
\begin{center}
    \textbf{\Large Registro delle decisioni}\\
    \vspace{10px}
    \begin{table}[h!]
        \centering
        \rowcolors{2}{logo!10}{logo!40}
        \renewcommand{\arraystretch}{1.8}
        \begin{tabular}{p{200px} p{200px}}
            \rowcolor{logo!70} \textbf{Codice} & \textbf{Decisione}                                                                                                                                                                           \\
            VE\_2021\_05\_16.1                 & Il proponente individua come apprezzabile e a maggior valore aggiunto l'eventuale scelta di passare ad un'architettura a microservizi. Il gruppo decide quindi di utilizzare i microservizi. \\
            VE\_2021\_05\_16.2                 & Il server-side rendering dovrà essere utilizzato il meno possibile, favorendo lo static generation.                                                                                          \\
            VE\_2021\_05\_16.3                 & Utilizzo di Chakra UI come framework utile al frontend.                                                                                                                                      \\
            VE\_2021\_05\_16.4                 & Utilizzo di ESLint con impostazioni Airbnb, sia nel frontend che nel backend, in quanto considerato migliore rispetto alle impostazioni ESLint che il gruppo sta utilizzando.                \\
            VE\_2021\_05\_16.5                 & Utilizzo di webhook per la gestione delle sessioni di pagamento.                                                                                                                             \\
            VE\_2021\_05\_16.6                 & Fissata una riunione per il 2021-05-21 in cui saranno presenti entrambi i rappresentanti dell'azienda proponente \textit{Red Babel}.
        \end{tabular}
    \end{table}
\end{center}

\end{document}